% Options for packages loaded elsewhere
\PassOptionsToPackage{unicode}{hyperref}
\PassOptionsToPackage{hyphens}{url}
\documentclass[
  a4paper,
]{ltjsarticle}
\usepackage{xcolor}
\usepackage[margin=25mm]{geometry}
\usepackage{amsmath,amssymb}
\setcounter{secnumdepth}{-\maxdimen} % remove section numbering
\usepackage{iftex}
\ifPDFTeX
  \usepackage[T1]{fontenc}
  \usepackage[utf8]{inputenc}
  \usepackage{textcomp} % provide euro and other symbols
\else % if luatex or xetex
  \usepackage{unicode-math} % this also loads fontspec
  \defaultfontfeatures{Scale=MatchLowercase}
  \defaultfontfeatures[\rmfamily]{Ligatures=TeX,Scale=1}
\fi
\usepackage{lmodern}
\ifPDFTeX\else
  % xetex/luatex font selection
\fi
% Use upquote if available, for straight quotes in verbatim environments
\IfFileExists{upquote.sty}{\usepackage{upquote}}{}
\IfFileExists{microtype.sty}{% use microtype if available
  \usepackage[]{microtype}
  \UseMicrotypeSet[protrusion]{basicmath} % disable protrusion for tt fonts
}{}
\makeatletter
\@ifundefined{KOMAClassName}{% if non-KOMA class
  \IfFileExists{parskip.sty}{%
    \usepackage{parskip}
  }{% else
    \setlength{\parindent}{0pt}
    \setlength{\parskip}{6pt plus 2pt minus 1pt}}
}{% if KOMA class
  \KOMAoptions{parskip=half}}
\makeatother
\usepackage{graphicx}
\makeatletter
\newsavebox\pandoc@box
\newcommand*\pandocbounded[1]{% scales image to fit in text height/width
  \sbox\pandoc@box{#1}%
  \Gscale@div\@tempa{\textheight}{\dimexpr\ht\pandoc@box+\dp\pandoc@box\relax}%
  \Gscale@div\@tempb{\linewidth}{\wd\pandoc@box}%
  \ifdim\@tempb\p@<\@tempa\p@\let\@tempa\@tempb\fi% select the smaller of both
  \ifdim\@tempa\p@<\p@\scalebox{\@tempa}{\usebox\pandoc@box}%
  \else\usebox{\pandoc@box}%
  \fi%
}
% Set default figure placement to htbp
\def\fps@figure{htbp}
\makeatother
\setlength{\emergencystretch}{3em} % prevent overfull lines
\providecommand{\tightlist}{%
  \setlength{\itemsep}{0pt}\setlength{\parskip}{0pt}}
% Glyph map for lualatex (newunicodechar). Maps symbols that may lack font glyphs.
\usepackage{newunicodechar}
\usepackage{amssymb}
\newunicodechar{∎}{\ensuremath{\blacksquare}}
\newunicodechar{✓}{\ensuremath{\checkmark}}
\newunicodechar{✗}{\ensuremath{\times}}
\newunicodechar{⟹}{\ensuremath{\Longrightarrow}}
\newunicodechar{⟺}{\ensuremath{\Longleftrightarrow}}
\newunicodechar{↪}{\ensuremath{\hookrightarrow}}
\newunicodechar{⌊}{\ensuremath{\lfloor}}
\newunicodechar{⌋}{\ensuremath{\rfloor}}
\newunicodechar{≃}{\ensuremath{\simeq}}
\newunicodechar{≈}{\ensuremath{\approx}}
\newunicodechar{≤}{\ensuremath{\leq}}
\newunicodechar{≥}{\ensuremath{\geq}}
\newunicodechar{≠}{\ensuremath{\neq}}
\newunicodechar{→}{\ensuremath{\rightarrow}}
\newunicodechar{←}{\ensuremath{\leftarrow}}
\newunicodechar{↔}{\ensuremath{\leftrightarrow}}
\newunicodechar{⊥}{\ensuremath{\perp}}
\newunicodechar{√}{\ensuremath{\surd}}
\newunicodechar{∑}{\ensuremath{\sum}}
\newunicodechar{∏}{\ensuremath{\prod}}
\newunicodechar{∞}{\ensuremath{\infty}}
\newunicodechar{∈}{\ensuremath{\in}}
\newunicodechar{∀}{\ensuremath{\forall}}
\newunicodechar{∣}{\ensuremath{\mid}}
\newunicodechar{γ}{\ensuremath{\gamma}}
\newunicodechar{λ}{\ensuremath{\lambda}}
\newunicodechar{μ}{\ensuremath{\mu}}
\newunicodechar{ρ}{\ensuremath{\rho}}
\newunicodechar{σ}{\ensuremath{\sigma}}
\newunicodechar{φ}{\ensuremath{\varphi}}
\newunicodechar{ψ}{\ensuremath{\psi}}
\newunicodechar{θ}{\ensuremath{\theta}}
\newunicodechar{Δ}{\ensuremath{\Delta}}
\newunicodechar{Π}{\ensuremath{\Pi}}
\newunicodechar{Λ}{\ensuremath{\Lambda}}
\newunicodechar{τ}{\ensuremath{\tau}}
\newunicodechar{ε}{\ensuremath{\varepsilon}}
\newunicodechar{ω}{\ensuremath{\omega}}
\newunicodechar{π}{\ensuremath{\pi}}
\newunicodechar{ν}{\ensuremath{\nu}}
\newunicodechar{±}{\ensuremath{\pm}}
\newunicodechar{×}{\ensuremath{\times}}
\newunicodechar{·}{\ensuremath{\cdot}}
\usepackage{bookmark}
\IfFileExists{xurl.sty}{\usepackage{xurl}}{} % add URL line breaks if available
\urlstyle{same}
\hypersetup{
  pdftitle={90度床からの自発的対称性の破れと縮退真空------微小シードによる対称性の破れと \textbackslash Delta=-2\textbackslash pi/N ロック（N=6{,}\textbackslash ldots{,}40 全数）},
  pdfauthor={Noriaki Kihara (WF System Co., Ltd.), ORCID 0009-0004-6753-4020},
  hidelinks,
  pdfcreator={LaTeX via pandoc}}

\title{90度床からの自発的対称性の破れと縮退真空------微小シードによる対称性の破れと
\(\Delta=-2\pi/N\) ロック（\(N=6,\ldots,40\) 全数）}
\author{Noriaki Kihara (WF System Co., Ltd.), ORCID 0009-0004-6753-4020}
\date{2026-09-12}

\begin{document}
\maketitle

\textbf{シリーズ:} 自己無撞着な関係波閉鎖系の基礎（論文5／全10本）\\
\textbf{著者:} 木原 範昭（WF System Co., Ltd.）　\textbf{日付:}
2026-09-12\\
\textbf{Version DOI:} 10.5281/zenodo.22729089\\
\textbf{Concept DOI:} 10.5281/zenodo.22729088\\
\textbf{ORCID:} 0009-0004-6753-4020\\
\textbf{Zenodo:} https://zenodo.org/records/22729089\\

\begin{center}\rule{0.5\linewidth}{0.5pt}\end{center}

\subsection{要旨}\label{ux8981ux65e8}

90度床は90°位相骨格を持つ自己無撞着な不安定相対平衡（論文1・2）である。同一床
\(Z_0\) に振幅 \(10^{-8}\) の\textbf{異なる向き}の微小シードを与えて厳密
one\_step
で走らせると、対称性が破れて指数的インフレーション（論文6）を経て、等振幅・位相分散の傾いた相対平衡（縮退真空の一つ）へロックする。\(N=6,\ldots,40\)
の全35本で、全シードがインフレーションを起こし、1step 回転角が厳密に
\(\Delta=-2\pi/N\)（誤差 \(\sim10^{-5}\)）・等振幅（比
\(\to1\)、\(\sim10^{-3}\)）の\textbf{同型平衡}へロックする一方、位相配置（真空の向き）はシード依存で相異する（真空差
中央 \(4.0^\circ\)--\(48.4^\circ\)、最大
\(76.9^\circ\)--\(180^\circ\)）。すなわち対称・不安定な床から、無限小シードが縮退した真空の一つを\textbf{自発的に選ぶ}。これは離散写像上の自発的対称性の破れ（メキシカンハット型）である。秩序変数
\(r=\sqrt{H_\perp/H}\) の実測ダイナミクスから Landau 型有効ポテンシャル
\(V(r)\)
を測定パラメータで再構成できる（\(N=6{:}\lambda=0.338,r_{\rm vac}=0.485\)／\(N=40{:}\lambda=0.076,r_{\rm vac}=0.330\)；アンサッツ）。500
step で未ロックの大 \(N\) も 2000 step で全数ロック（\(N=22\) は step
\(663\)、\(N=40\) は step
\(1299\)）。分類：新仮説の実測支持（SSB像・全数確認）。

\begin{center}\rule{0.5\linewidth}{0.5pt}\end{center}

\subsection{1. 課題}\label{ux8ab2ux984c}

論文1で90度床は「振幅駆動 \(\sin(2\Delta\varphi)\)
が全項個別に消える不安定な相対平衡（90°位相骨格）」、論文2で「点火の資格＝不安定な相対平衡」と確定した。本論文は、その不安定性が具体的に何を生むか------\textbf{無限小シードが何を選ぶか}------を全
\(N\) で調べ、メキシカンハット型 SSB として読む。

\subsection{2.
力学とシード（制御実験）}\label{ux529bux5b66ux3068ux30b7ux30fcux30c9ux5236ux5fa1ux5b9fux9a13}

写像は正本 one\_step（論文1 §2、SHA \texttt{1abf2353…}）。床 \(Z_0\) は
\texttt{full\_N3\_N40\_sweep} の step0（make\_parent 固有モード床、論文2
§6）。制御変数は初期シードのみ（物理式無変更、判定規則は力学に入れず走行後分析）：各
\(N\) について \(Z_0+10^{-8}(\xi+i\eta)\)（\(\xi,\eta\)
乱数、規格化）を向きの異なる5通り与え、適応停止（等振幅かつゲージ不変な相対位相変化率が閾値以下、上限
2600 step）まで走らせる。各シードの終状態で
lock\_step・\(H_\perp/H\)・振幅比・\(\Delta\)（1step
回転角）・\(|\Delta+2\pi/N|\) と、相対位相配置のシード間差を測る。

\subsection{3. 機構------メキシカンハット型
SSB}\label{ux6a5fux69cbux30e1ux30adux30b7ux30abux30f3ux30cfux30c3ux30c8ux578b-ssb}

90度床＝90°骨格の不安定な相対平衡（模式的にはソンブレロの頂上）、谷＝縮退した相対平衡の族（等振幅・位相分散・傾いた平衡。共通不変量
\(\Delta=-2\pi/N\)・等振幅を持つ）。無限小シードが対称性を破って系を不安定固有方向へ転がし（インフレーション、論文6）、谷の一点へロックする------不安定な対称極大から縮退した下位平衡の一つを自発的に選ぶ。谷（縮退真空多様体）は有限離散な写像対称性軌道ではなく\textbf{傾いた相対平衡の連続モジュライ}であることを
\(N=6\) で確認した（5終状態は \(6!=720\)
頂点置換＋大域位相＋複素共役のどれでも互いに写らない：\texttt{頂点置換による縮退真空同定\_N6\_20260913/}）。これはメキシカンハットの谷が本来
連続縮退多様体であることと整合する（谷が単一の離散対称性軌道をなすという主張ではない）。秩序変数
\(r=\sqrt{H_\perp/H}\)（\(H_\perp\)＝床平面に直交する成分のエネルギー）の実測から
Landau 型

\[V(r)=-\tfrac12\lambda r^2+\frac{\lambda}{4\,r_{\rm vac}^2}r^4,\qquad \lambda=\ln|\mu_{\max}|,\ r_{\rm vac}=\sqrt{H_\perp/H|_{\rm sat}}\]

を再構成できる（\(N=6{:}\lambda=0.338,r_{\rm vac}=0.485\)／\(N=40{:}\lambda=0.076,r_{\rm vac}=0.330\)）。これは測定量で係数を同定した\textbf{有効記述（アンサッツ）}である。\(V(r)\)
の第一原理導出は今後の検討課題とする。\(r=\sqrt{H_\perp/H}\)
は床からの離脱量を表す一つのスカラー座標であり、図のメキシカンハットは
SSB 構図の\textbf{模式的可視化}であって、\(r\)
が全縮退真空に共通する厳密な動径座標であることを意味しない（実際
\(H_\perp/H\) は真空ごとに \(0.19\)--\(0.38\)
と相異する、§7）。ここで「真空」は有効ポテンシャルの縮退極小を指す数学的呼称で、物理的真空の特徴づけは今後の検討課題とする。なお
\(\lambda=\ln|\mu_{\max}|\)（\(\mu_{\max}\)＝床ヤコビアンのスペクトル半径、\(iK\)
の \(\mu\)
とは別量）は論文6のインフレ増幅率そのもので、\textbf{SSBポテンシャルの曲率＝インフレの増幅率}が同一の
\(\mu_{\max}\) で結ばれる。

\begin{figure}
\centering
\pandocbounded{\includegraphics[keepaspectratio,alt={図1: メキシカンハット型SSB（実測再構成, N=6）。左=V(r)上を実測秩序変数が床（不安定頂点）から谷へ転落、右=ソンブレロ面と縮退真空リング（シード5点）。}]{figs/p5_ssb_degenerate_vacuum_ja_fig1.png}}
\caption{図1: メキシカンハット型SSB（実測再構成,
N=6）。左=V(r)上を実測秩序変数が床（不安定頂点）から谷へ転落、右=ソンブレロ面と縮退真空リング（シード5点）。}
\end{figure}

\begin{figure}
\centering
\pandocbounded{\includegraphics[keepaspectratio,alt={図2: 同 N=40（λ=0.076, r\_vac=0.330）。}]{figs/p5_ssb_degenerate_vacuum_ja_fig2.png}}
\caption{図2: 同 N=40（λ=0.076, r\_vac=0.330）。}
\end{figure}

\textbf{図1・2} メキシカンハット型
SSB（実測再構成、\(N=6\)／\(N=40\)）。

\subsection{\texorpdfstring{4. 決定的検証（同一床＋異なるシード →
縮退真空、\(N=6\)）}{4. 決定的検証（同一床＋異なるシード → 縮退真空、N=6）}}\label{ux6c7aux5b9aux7684ux691cux8a3cux540cux4e00ux5e8aux7570ux306aux308bux30b7ux30fcux30c9-ux7e2eux9000ux771fux7a7an6}

\(N=6\) 床に振幅 \(10^{-8}\)
の向きの違うシードを与えた終状態：全シードが
\(\Delta=-1.047198=-2\pi/6\)（\(|\Delta+2\pi/N|\le4.3\times10^{-7}\)）・振幅比
\(1.0000\)--\(1.0006\)
の\textbf{同型平衡}へロック。一方、真空の向き（相対位相配置）のシード間差は
中央 \(1.6^\circ\)（seed1 vs 0）から \(112.7^\circ\)（seed5 vs
0）まで\textbf{相異する}。不変量（\(\Delta=-2\pi/N\)・等振幅）は真空共通、破れた自由度（位相配置の向き）がシードで選ばれる。

\begin{figure}
\centering
\pandocbounded{\includegraphics[keepaspectratio,alt={図3: N=6 同一床＋異なる微小シード5本。左=各シードのロック配置（剛体回転除去、等半径だが位相配置相異＝縮退真空）、中=全シードΔ=-2π/N、右=全N6-40の真空差のN依存。}]{figs/p5_ssb_degenerate_vacuum_ja_fig3.png}}
\caption{図3: N=6
同一床＋異なる微小シード5本。左=各シードのロック配置（剛体回転除去、等半径だが位相配置相異＝縮退真空）、中=全シードΔ=-2π/N、右=全N6-40の真空差のN依存。}
\end{figure}

\textbf{図3} SSB 縮退真空（実データ、\(N=6\)）。

\subsection{\texorpdfstring{5.
全数確認（\(N=6,\ldots,40\)、35/35）}{5. 全数確認（N=6,\textbackslash ldots,40、35/35）}}\label{ux5168ux6570ux78baux8a8dn6ldots403535}

\(N=6\)--\(40\) の全35本（各 \(N\) 向きの異なる5シード＝計
\(35\times5=175\) 走行）で、SSB の3条件がすべて成立（バッチ集計
\texttt{ssb\_batch\_summary.csv}）：

\begin{itemize}
\tightlist
\item
  \textbf{全シードがインフレーションを起こす}（床から明確に離脱、\texttt{onset\_ok=True}
  全35）。
\item
  \textbf{全て等振幅}（振幅比 \(<1.01\)、\texttt{all\_equal\_amp=True}
  全35）・\textbf{1step 回転角が厳密に
  \(\Delta=-2\pi/N\)}（\texttt{all\_lock\_2piN=True}
  全35、\(|\Delta+2\pi/N|<10^{-3}\)）。
\item
  \textbf{位相配置（真空の向き）はシード依存で相異}：真空差 中央
  \(4.0^\circ\)（\(N=40\)）--\(48.4^\circ\)（\(N=10\)）、最大
  \(76.9^\circ\)（\(N=26\)）--\(180.0^\circ\)。
\end{itemize}

すなわち不安定な床から、写像の対称性が許す縮退真空の一つを無限小シードが自発的に選ぶ（35/35
で SSB 確認）。

\begin{figure}
\centering
\pandocbounded{\includegraphics[keepaspectratio,alt={図4: 全N6-40のSSBバッチ。左=真空差(中央/最大)のN依存、右=実測Δ\_tailが理論-2π/Nに全N一致(誤差\textasciitilde1e-5)・振幅比→1。}]{figs/p5_ssb_degenerate_vacuum_ja_fig4.png}}
\caption{図4:
全N6-40のSSBバッチ。左=真空差(中央/最大)のN依存、右=実測Δ\_tailが理論-2π/Nに全N一致(誤差\textasciitilde1e-5)・振幅比→1。}
\end{figure}

\textbf{図4} 全 \(N=6\)--\(40\) の SSB（35/35）。

\subsection{\texorpdfstring{6. 長時間検証（\(N=22,40\)・2000
step）}{6. 長時間検証（N=22,40・2000 step）}}\label{ux9577ux6642ux9593ux691cux8a3cn22402000-step}

500 step で未ロックだった大
\(N\)（\(N=22,30,\ldots,40\)）は打ち切りが原因で、2000 step
まで無停止走行すると全数ロックする。\(N=22\) は step \(663\)、\(N=40\)
は step \(1299\) で完全ロック（\(\Delta_{\rm tail}\) 厳密
\(-2\pi/N\)、等振幅比 \(1.0000\)、床/末尾の剛体回転残差
機械ゼロ）。対照テストで初期値が旧走行と step0 bit 一致（step1 差
\(\sim10^{-16}\)、以後インフレが指数増幅）を確認し忠実再現を担保した。

\subsection{7.
主張分類・未決}\label{ux4e3bux5f35ux5206ux985eux672aux6c7a}

\begin{itemize}
\tightlist
\item
  分類：新仮説の実測支持（SSB像・全数確認）。判定：保持。
\item
  断定：本系は自発的対称性の破れを示す（構図------90°骨格の不安定な相対平衡→縮退した下位平衡の族→微小シードによる自発選択→全真空同型
  \(\Delta=-2\pi/N\)・等振幅------が機械精度で成立、35/35。メキシカンハットは模式的有効模型、§3）。\textbf{要点は、同じ90°床から終端面（真空の向き）がシード依存で一意に定まらないこと。}
\item
  真空多様体（谷）の構造：\(N=6\) の5終状態は \(6!=720\)
  頂点置換＋大域位相（＋複素共役）のどれでも互いに写らない（最小残差
  \(0.05\)--\(0.83\)）。ゆえ谷は有限離散な写像対称性軌道ではなく\textbf{傾いた相対平衡の連続モジュライ}（\(H_\perp/H=0.19\)--\(0.38\)
  が連続的）。これはメキシカンハットの谷が本来
  連続縮退多様体であることと整合し、SSB
  を否定しない。モジュライの次元・計量の同定は今後の検討課題。
\item
  対称群の詳細な分類は本論文の範囲外とする。不変量は \(\Delta=-2\pi/N\)
  と等振幅で、シード依存で変わるのは真空の向き（傾き \(H_\perp/H\) が
  \(0.19\)--\(0.38\)、\(N=6\)）。「真空」は有効ポテンシャルの縮退極小の数学的呼称で、物理的真空は未特徴づけ。\(V(r)\)
  はアンサッツ。
\end{itemize}

\subsection{関連研究（構造的一致・導出元でない）}\label{ux95a2ux9023ux7814ux7a76ux69cbux9020ux7684ux4e00ux81f4ux5c0eux51faux5143ux3067ux306aux3044}

\begin{itemize}
\tightlist
\item
  自発的対称性の破れ／Landau 理論／Goldstone・Higgs（Landau 1937;
  Goldstone 1961; Anderson 1963; Higgs
  1964）：本節は離散写像上のSSBの一例で、再構成した有効ポテンシャルは
  Ginzburg--Landau 型ソンブレロ。
\item
  蔵本モデル（同期）：ロック＝同期。
\end{itemize}

\subsection{参考文献}\label{ux53c2ux8003ux6587ux732e}

\begin{enumerate}
\def\labelenumi{\arabic{enumi}.}
\tightlist
\item
  木原範昭「自己無撞着な関係波閉鎖系におけるインフレーション的急拡大の機構」Concept
  DOI 10.5281/zenodo.22112008.
\item
  L. D. Landau, ``On the theory of phase transitions'', 1937; J.
  Goldstone, Nuovo Cimento 19 (1961) 154; P. W. Anderson, Phys. Rev.~130
  (1963) 439; P. W. Higgs, Phys. Rev.~Lett. 13 (1964) 508.
\item
  Y. Kuramoto, \emph{Chemical Oscillations, Waves, and Turbulence},
  Springer (1984).
\end{enumerate}

\subsection{再現}\label{ux518dux73fe}

検証スクリプト・図・集計は
\texttt{位相ロック全N確認\_相対平衡\_20260912/自発的対称性の破れ\_シード依存縮退真空\_20260912/}（\texttt{verify\_ssb\_seed\_vacua\_20260912.py}、\texttt{make\_figure\_ssb.py}、\texttt{make\_figure\_mexican\_hat.py}/\texttt{\_N40.py}、\texttt{REPORT.md}、\texttt{SHA256SUMS.txt}、\texttt{run\_all.sh}）と、その配下の
\texttt{全N6\_40\_SSBバッチ\_20260912/}（\texttt{run\_ssb\_batch\_N6\_N40\_20260912.py}、\texttt{per\_N/N\{6..40\}.txt}、\texttt{ssb\_batch\_summary.csv}、\texttt{fig\_ssb\_batch\_allN.png}）。長時間は
\texttt{N22\_2000step検証走行\_20260912/}・\texttt{N40\_2000step検証走行\_20260912/}。真空多様体が連続モジュライである（有限離散対称性軌道でない）ことの
\(N=6\) 決定的テストは
\texttt{頂点置換による縮退真空同定\_N6\_20260913/}（\texttt{verify\_vacua\_permutation\_orbit\_N6\_20260913.py}＝正本
one\_step を SHA 照合 import・\(S_6\) 総当たりの読出しのみ）。床データは
make\_parent 床（論文2 §6）。走行は正本
\texttt{run\_N3\_N40\_stage123\_v1.py}（SHA256
\texttt{1abf2353fee2e4f56f05e7a6f149fd086885136beb61ab571b48a56b09691567}）と同一・初期シードのみ制御。

\end{document}
